% ==============================================================================
% BỘ CÔNG THỨC HÓA HỌC DẠNG KHUNG (SKELETAL) & ION LƯỠNG CỰC (ZWITTERION)
% Bi soạn bởi Thầy TRẦN VĂN THẠNH - 0777.470.803
% Sử dụng gói chemfig trong LaTeX
% ==============================================================================

\documentclass[12pt,a4paper]{article}
\usepackage[utf8]{inputenc}
\usepackage[vietnamese]{babel}
\usepackage{amsmath,amssymb}
\usepackage{chemfig}
\usepackage[scaled=0.95]{helvet}
\usepackage{geometry}
\geometry{top=2cm,bottom=2cm,left=2cm,right=2cm}

% Cấu hình ChemFig chuẩn sách giáo khoa mới (Kết nối tri thức / Cánh diều / Chân trời sáng tạo)
\setchemfig{
    atom sep=2.2em,           % Khoảng cách giữa các nguyên tử (độ dài liên kết)
    bond offset=1.8pt,         % Khoảng cách hở giữa liên kết và chữ
    bond style={line width=1.0pt, line cap=round, line join=round}, % Độ dày nét vẽ 1.0pt
    double bond sep=2.5pt      % Khoảng cách giữa 2 nét liên kết đôi
}
% Đổi font chữ các nguyên tử sang Sans-Serif chuẩn
\renewcommand*\printatom[1]{\ensuremath{\mathsf{#1}}}

\begin{document}

\section*{CÁC CÔNG THỨC HÓA HỌC TRONG ĐỀ THI}

\subsection*{1. Từng công thức riêng lẻ}

\begin{itemize}
    \item \textbf{Phương án A (Zwitterion của $\beta$-alanine / Axit 3-aminopropanoic):}
    \begin{center}
        \chemfig{^{+}H_3N-[:30]-[:-30]-[:30]COO^{-}}
    \end{center}
    Mã ChemFig:
    \begin{verbatim}
    \chemfig{^{+}H_3N-[:30]-[:-30]-[:30]COO^{-}}
    \end{verbatim}

    \item \textbf{Phương án B (Zwitterion của Alanine / Axit 2-aminopropanoic - Dạng chuẩn Hóa học):}
    \begin{center}
        \chemfig{-[:-30](-[:-90]NH_3^{+})-[:30]COO^{-}}
    \end{center}
    Mã ChemFig:
    \begin{verbatim}
    \chemfig{-[:-30](-[:-90]NH_3^{+})-[:30]COO^{-}}
    \end{verbatim}

    \item \textbf{Phương án B (Phiên bản theo đúng ảnh đề thi in thiếu số 3: $NH^+$):}
    \begin{center}
        \chemfig{-[:-30](-[:-90]NH^{+})-[:30]COO^{-}}
    \end{center}
    Mã ChemFig:
    \begin{verbatim}
    \chemfig{-[:-30](-[:-90]NH^{+})-[:30]COO^{-}}
    \end{verbatim}

    \item \textbf{Phương án C (Zwitterion của Axit 2,3-diaminopropanoic):}
    \begin{center}
        \chemfig{^{+}H_3N-[:30](-[:90]NH_2)-[:-30]-[:30]COO^{-}}
    \end{center}
    Mã ChemFig:
    \begin{verbatim}
    \chemfig{^{+}H_3N-[:30](-[:90]NH_2)-[:-30]-[:30]COO^{-}}
    \end{verbatim}

    \item \textbf{Phương án D (Aminomethanol - $H_2N-CH_2-OH$):}
    \begin{center}
        \chemfig{H_2N-[:30]-[:-30]OH}
    \end{center}
    Mã ChemFig:
    \begin{verbatim}
    \chemfig{H_2N-[:30]-[:-30]OH}
    \end{verbatim}
\end{itemize}

\subsection*{2. Bố trí 4 phương án trên 1 hàng ngang (Chuẩn đề thi trắc nghiệm)}

\noindent
\textbf{A.} $\vcenter{\hbox{\chemfig{^{+}H_3N-[:30]-[:-30]-[:30]COO^{-}}}}$ \qquad\quad
\textbf{B.} $\vcenter{\hbox{\chemfig{-[:-30](-[:-90]NH_3^{+})-[:30]COO^{-}}}}$ \qquad\quad
\textbf{C.} $\vcenter{\hbox{\chemfig{^{+}H_3N-[:30](-[:90]NH_2)-[:-30]-[:30]COO^{-}}}}$ \qquad\quad
\textbf{D.} $\vcenter{\hbox{\chemfig{H_2N-[:30]-[:-30]OH}}}$

\subsection*{3. Bố trí lưới 2x2 (Khớp với bố cục đề gốc)}

\begin{center}
\begin{tabular}{ll}
\textbf{A.} $\vcenter{\hbox{\chemfig{^{+}H_3N-[:30]-[:-30]-[:30]COO^{-}}}}$ &
\qquad\qquad \textbf{B.} $\vcenter{\hbox{\chemfig{-[:-30](-[:-90]NH_3^{+})-[:30]COO^{-}}}}$ \\[3.5em]
\textbf{C.} $\vcenter{\hbox{\chemfig{^{+}H_3N-[:30](-[:90]NH_2)-[:-30]-[:30]COO^{-}}}}$ &
\qquad\qquad \textbf{D.} $\vcenter{\hbox{\chemfig{H_2N-[:30]-[:-30]OH}}}$
\end{tabular}
\end{center}

\end{document}
